\usepackage[%
            a4paper,
            width=150mm,
            left=35mm,
            bottom=30mm,
            top=20mm,
            %headheight=10mm,
            headsep=25mm,
            includehead
            ]{geometry}
\usepackage[utf8]{inputenc}                 %% Standardinputformat in UTF-8
\usepackage[german]{babel}                  %% Umstellung der Standardisierten Titel auf Deutsch
\usepackage[dvipsname]{xcolor}
\usepackage{multirow}
\usepackage{lastpage}
\usepackage{graphicx}
\usepackage{epstopdf}
\usepackage{pdfpages}
%\usepackage{enumitem}
\usepackage{ragged2e}
\usepackage{changepage}
\usepackage{titling}
\usepackage{tabularx}
\usepackage{booktabs}
\usepackage{longtable}
\usepackage{enumerate}                      %% Feinere Nummerierungen
\usepackage{textpos}
\usepackage{fancyhdr}                       %% Ermöglicht die feine Steuerung der Kopf- und Fusszeilen
\pagestyle{fancy}
\fancyhead[L,C]{}
\fancyhead[R]{%
\vspace{-29mm}
\tiny
\begin{tabular}{m{6mm}m{29mm}}
  \multirow{4}{*}{\includegraphics[width=6mm]{layout/schild.eps}}&Kantonsschule
  Büelrain\\
  &\thetitle\\
  &\thedate\\
  &\thepage/\pageref{LastPage}\\
\end{tabular}
\normalsize
}
\fancyfoot[L,C,R]{}

\fancypagestyle{firstpage}{%
\lhead{%
    \begin{tikzpicture}[remember picture, overlay]
      \node[draw=none] at (-2,0) (leu)
      {\includegraphics[width=2.25cm]{layout/leu.eps}\raisebox{1.3cm}{\includegraphics[width=6mm]{layout/schild.eps}}};
      \node[draw=none, anchor=west, right= -3mm of leu] (absender) {%
      \footnotesize
      \color{black}
      \begin{tabular}{l}
        Kanton Zürich\\
        \large \textbf{Kantonsschule Büelrain}\\
        Fachschaft WR\\
      \end{tabular}
      };
      \node[draw=none, anchor=west, below right = .25cm and -5.5cm of absender] (autor) {%
      \footnotesize
      \color{black}
      \begin{tabular}{l}
        \theauthor\\
        Rosenstrasse 1\\
        8400 Winterthur\\
        jacques.mock@kbw.ch\\
        www.kbw.ch\\
      \end{tabular}
      };
    \end{tikzpicture}
    }
\rhead{}
}

\fancypagestyle{inhaltsverzeichnis}{%
            \lhead{}
            \rhead{%
            \tiny
            \begin{tabular}{m{6mm}m{2.9cm}}
                \includegraphics[width=6mm]{layout/schild.eps}&\thepage/\pageref{LastPage}\\
            \end{tabular}
            \normalsize
              }
            }
\renewcommand{\headrulewidth}{0pt}

\usepackage{helvet}                         %% Schriftart Helvetica
\renewcommand{\familydefault}{\sfdefault}   %% Definition Helvetica als Standardschrift
%\usepackage{a4}                             %% Seitenformat A4
\usepackage[T1]{fontenc}                    %% Westeuropäischer Zeichensatz
%

\usepackage[backend=biber,                  %% Steuerung der Darstellung der Fussnoten und des Literaturverzeichnisses
            isbn=false,
            autopunct=false,
            sorting=none,
            style=verbose-ibid]{biblatex}
\DeclareNameAlias{author}{last-first}        %% Voranstellung des Nachnamens für die Sortierung in Verzeichnissen
\bibliography{references}                    %% Deklaration der Literaturdatenbank

\usepackage{tikz}                            %% Graphikprogramm
\usetikzlibrary{positioning,
                fit,
                calc,
                shapes,
                shadows,
                arrows,
                trees,
                decorations.markings,
                decorations.pathreplacing,
                decorations.text}
\usepackage{pgfplots}                

\usepackage{microtype}                      %% Schönere Darstellung der Buchstaben

\usepackage[%ps2pdf                         %% Steuerung der PDF-Ausgabe
	]{thumbpdf}
\usepackage[%                               %% Inhaltsverzeichnis als Hyperlink
	%ps2pdf,
	pdfa]{hyperref}
	%pagebackref]{hyperref}
\hypersetup{%                               %% Darstellung der Links
   colorlinks=true,        % Einf\"arbung von Verkn\"upfungen
   linkcolor=black,        % stzt die Linkfarbe auf schwarz
   citecolor=black,        % Literaturzitate in schwarz
   urlcolor=black,         % URL werden schwarz dargestellt
   breaklinks=true,        % erm\"oglicht mehrzeilige Links
   pdfpagemode=UseThumbs,  % Anzeige der Piktogramme
   bookmarksopen=true,     % Anzeige aller Ebenen
   bookmarksnumbered=true, % Anzeige der Abschnittsnummern
   pdfstartpage={1},       % Startseite
   pdftitle={Vorlage Rechtstext PVBR 2bH},
   pdfsubject={Prüfung},
   pdfauthor={Jacques Mock Schindler},%
   pdfkeywords={},
   pdfcreator={Adobe-Acrobat-Distiller},%
   pdfproducer={LaTeX with hyperref and thumbpdf}
}
\usepackage{lastpage}                        %% von Seiten

\renewcommand\labelitemi{--}
\sloppy
\usepackage{xcolor}
\definecolor{zhblue}{cmyk}{100,0,0,0}
